\section{Introduction}
A Monte Carlo simulation produces numbers, and it is tempting to read them as measurements. A curve with a peak is read as a transition, and the temperature of the peak as the transition temperature. That reading needs two things to be true. The Markov chain has to sample the distribution we intend, and the observable we extract has to identify the physical quantity we intend. A simulation does not become trustworthy because it runs, equilibrates and produces a plausible-looking figure.

We check both conditions on a system where exact references exist for each. The system is the two-dimensional Ising lattice gas at fixed composition, evolved by pair-exchange Metropolis moves that conserve the number of each type. It is also a minimal statistical-physics model of residential segregation, in which a ``social temperature'' sets the balance between preference for like neighbours and tolerance of mixing \citep{schelling1971,vinkovic2006,stauffer2007,dallasta2008,gauvin2009}. In that literature a transition or demixing temperature is often read off simulation output. For this model we can say what the right answers are: the canonical law can be computed exactly on small tori, and the fixed-composition coexistence temperature is known in closed form from the work of Onsager and Yang \citep{onsager1944,yang1952}. That makes it possible to ask how far a standard simulation procedure is from them, and why.

\paragraph{Question 1: does the sampler target the canonical law?}
Exchanging two unlike spins is the same as flipping both. If the two sites are not neighbours the energy change is the sum of the two single-flip changes. If they are neighbours, the bond between them is unchanged by the exchange but is counted twice by that sum, and the exact change is the sum plus $4J$ (\cref{prop:delta}). The correction touches only a fraction $O(N^{-2})$ of uniformly random pair proposals, so it is small, and the interesting question is whether it matters statistically. A one-step perturbation of order $N^{-2}$ does not by itself imply a stationary-distribution perturbation of the same order; that depends on how fast the chain mixes. We therefore answer by proof where we can (detailed balance, irreducibility, a bound on the one-step perturbation), by exact stationary laws on tori small enough to solve, and by replicated simulation at production sizes. This is the foundation for the second question, whose conclusions are only meaningful if the chain is correct.

\paragraph{Question 2: given a correct sampler, does the variance peak locate the transition?}
A common estimate of a transition temperature is the location of the maximum of the energy variance over a temperature grid. At $f=1/2$ this is tied to the critical point. At other compositions the relevant reference is the fixed-composition coexistence (binodal) temperature $\Tb(f)$, and there is no divergence to locate: a finite system at fixed composition passes from a homogeneous state to coexistence through droplet formation, and this is rounded and displaced by finite-size effects \citep{binder1987,biskup2002}. So whether the variance peak tracks $\Tb(f)$ is a finite-size question rather than an identity. The exact curve is also nearly flat, with $\Tb(0.2)=\EbFifth$ against $\Tc=\EbHalf$, which means that a composition dependence of a few tenths in a simulated peak inside $0.2\le f\le0.8$ cannot be a property of the transition. Settling this requires the exact reference, replicate runs and a comparison with the resolution of the estimator, and it requires knowing whether the sampling window is long enough: the energy autocorrelation time can exceed the 300-sweep window of the protocol we study.

\paragraph{What we find.} The corrected exchange kernel has the canonical law as its unique stationary distribution, and we prove it. The version without the adjacency term has a stationary law far from canonical on small tori, but the effect on mean energy fades with lattice size under random-pair proposals and we cannot detect it for $N\ge32$ with the replication we could afford. The variance peak is noisy, its single-run scatter is larger than the entire composition dependence of the exact curve in the central range, and at dilute compositions it sits systematically below the exact curve. The long equilibrium reference does not remove that offset.

\paragraph{Contributions.} We do not propose a new model or algorithm. Our contributions are the following:
\begin{enumerate}[leftmargin=2em]
\item \textbf{Exact exchange energetics and detailed balance.} We derive the exact exchange energy (\cref{prop:delta}) and its attainable values (\cref{cor:values}), and prove that the exchange kernel is reversible, irreducible and aperiodic with the canonical law as unique stationary distribution (\cref{thm:correct}). Verified exhaustively on all $\EoneConfigs$ local configurations of a $6\times6$ torus with no mismatch.
\item \textbf{Local kernel error versus stationary law.} Without the adjacency term, detailed balance fails by exactly $e^{\pm4\beta J}$ (\cref{prop:flux}) and the one-step perturbation is at most $4N^{-2}(1-e^{-4\beta J})$ (\cref{prop:footprint}). We separate this local statement from the stationary one: exact stationary laws on tori of up to $8008$ states are $\EtwoPairMin$--$\EtwoPairMax$ (random-pair) and $\EtwoNNMin$--$\EtwoNNMax$ (nearest-neighbour proposals) from canonical in total variation, against $\EtwoMaxTVcorrect$ for the corrected kernel; replicated simulation gives a mean-energy shift of $5.5$--$6.3\%$ at $N=8$ and $1.0$--$2.2\%$ at $N=16$, and none we can detect at $N\ge32$.
\item \textbf{The variance peak against the exact binodal.} Over $\EfiveR$ replicates, three sizes and nine compositions, the single-run standard deviation is $\EfiveSdMin$--$\EfiveSdMax$; the central-range location is within $\EfiveCenBiasMax$ of the exact curve but too coarse to resolve its composition dependence; at dilute compositions it lies $0.17$--$0.29$ below it, with $|z|\ge\EfiveExtZMin$ in all nine cells.
\item \textbf{Equilibration and finite-size effects.} In long traces the 300-sweep window is shorter than the energy autocorrelation time in the worst cells and biases the variance. An equilibrium reference run shows the dilute offset is not a protocol artefact; it is consistent with finite-size droplet formation, though we do not test the scaling (\cref{sec:results-equil}).
\item \textbf{A reproducible pipeline.} Code, raw output and a generator for every table and figure (\cref{app:repro}).
\end{enumerate}
Contributions 1 and 2 give the computational foundation and quantify how a local kernel error propagates to equilibrium; contributions 3 and 4 ask whether the standard estimator locates the transition.

\paragraph{Organisation.} \Cref{sec:related} gives background. \Cref{sec:theory} treats the exchange kernel and \cref{sec:thermo} the thermodynamic reference and the estimator. \Cref{sec:protocol} describes the experiments. \Cref{sec:results} examines the sampler, \cref{sec:results-demix} the estimator, and \cref{sec:discussion} states what the study does and does not establish.
